% Source: https://artofproblemsolving.com/community/c3h1280937_2005_state_team_4
% Stable identity: aops-topic-1280937-q01
Let the ages be $A$, $J$, $M=A+J$. Since $\dfrac {J}{2}+1$ is a prime (an integer), $J$ is even and $\dfrac {J}{2}\in\{1,2,4,6,10,\dots\}$, i.e. $J\in\{2,4,8,12,20,\dots\}$. John is not prime and not a teenager, and $J<A\le19$, so $J\in\{4,8,12\}$. Angie is a prime teenager: $A\in\{13,17,19\}$.

Requiring $A-J$ prime leaves the pairs $(A,J)=(17,4)$, $(13,8)$, $(19,8)$, $(17,12)$, $(19,12)$. Now check $M=A+J$ (not prime) and $J+M=A+2J$ (prime):
\begin{itemize}
\item $(17,4)$: $A+2J=25$, not prime.
\item $(13,8)$: $M=21$, not prime and not a teenager; $A+2J=29$, prime. Works.
\item $(19,8)$: $A+2J=35$, not prime.
\item $(17,12)$ and $(19,12)$: $M=29$ or $31$, which are prime.
\end{itemize}
So John is $\boxed{8}$ (Angie $13$, Michael $21$).
